% Part: many-valued-logic
% Chapter: three-valued-logics
% Section: multiple-designation

\documentclass[../../../include/open-logic-section]{subfiles}

\begin{document}

\olfileid[fr]{mvl}{thr}{mul}

\olsection{Désigner d'autres valeurs que~\texorpdfstring{$\True$}{𝕋}}

Dans toutes les logiques rencontrées
jusqu'ici, l'ensemble des valeurs de
vérité désignées était
$V^+=\{\True\}$ : un énoncé comptait
comme vrai si et seulement si sa
valeur de vérité était~$\True$.
Mais on pourrait aussi compter un
énoncé comme vrai dès lors que sa
valeur n'est pas~$\False$. On peut
ainsi obtenir une logique en posant,
par exemple, $V^+=\{\True,\Undef\}$
dans sa matrice.

\begin{defn}
La \emph{logique du paradoxe}~$\LogLP$
est définie par la matrice suivante :
\begin{enumerate}
  \item Le langage propositionnel usuel
  $\Lang L_0$ avec $\lnot$, $\land$,
  $\lor$ et $\lif$.
  \item L'ensemble des valeurs de vérité
  $V=\{\True,\Undef,\False\}$.
  \item $\True$ et $\Undef$ sont
  désignées, c'est-à-dire
  $V^+=\{\True,\Undef\}$.
  \item Les fonctions de vérité sont
  celles de la logique forte de Kleene.
\end{enumerate}
\end{defn}

\begin{defn}
La \emph{logique du non-sens}
de Halld\'en~$\LogHal$ est définie
par la matrice suivante :
\begin{enumerate}
  \item Le langage propositionnel usuel
  $\Lang L_0$ avec $\lnot$, $\land$,
  $\lor$, $\lif$ et un connecteur
  unaire~$+$.
  \item L'ensemble des valeurs de vérité
  $V=\{\True,\Undef,\False\}$.
  \item $\True$ et $\Undef$ sont
  désignées, c'est-à-dire
  $V^+=\{\True,\Undef\}$.
  \item Les fonctions de vérité sont
  celles de la logique faible de
  Kleene, auxquelles s'ajoute
  l'opérateur «~est dépourvu de sens~» :
  \begin{center}
    \begin{tabular}{c|c}
      $\tf{+}$ & \\
      \hline
      $\True$ & $\False$ \\
      $\Undef$ & $\True$ \\
      $\False$ & $\False$
    \end{tabular}
  \end{center}
\end{enumerate}
\end{defn}

Contrairement aux logiques de Kleene
dont elles reprennent les tables de
vérité, ces logiques \emph{ont} des
tautologies.

\begin{prop}\ollabel{prop:LP-taut-CL}
Les tautologies de~$\LogLP$ sont
exactement celles de la logique
propositionnelle classique, sur leur
langage commun.
\end{prop}

\begin{proof}
  D'après la \olref[syn][sub]{prop:mvl-cl},
  si $\Entails[\LogLP] !A$, alors
  $\Entails[\LogCL] !A$. Pour établir
  la réciproque, montrons que, s'il
  existe une valuation
  $\pAssign{v}\colon\PVar\to
  \{\False,\True,\Undef\}$ telle que
  $\pValue{v}(!A)[\LogKs]=\False$,
  il existe une valuation classique
  $\pAssign{v'}\colon\PVar\to
  \{\False,\True\}$ telle que
  $\pValue{v'}(!A)[\LogCL]=\False$.
  Cela suffit pour~$\LogLP$ : elle a
  les mêmes fonctions de vérité que
  $\LogKs$, et $\False$ est sa seule
  valeur non désignée. En effet, si
  $\Entails/[\LogLP] !A$, une valuation
  $\pAssign{v}$ donne
  $\pValue{v}(!A)[\LogLP]
  =\pValue{v}(!A)[\LogKs]=\False$.
  L'affirmation à démontrer fournit
  alors une valuation classique
  $\pAssign{v'}$ pour laquelle
  $\pValue{v'}(!A)[\LogCL]=\False$;
  par conséquent,
  $\Entails/[\LogCL] !A$.

  Pour prouver cette affirmation,
  définissons d'abord $\pAssign{v'}$
  par
  \[
    \pAssign{v'}(p)=
    \begin{cases}
      \True & \text{si }\pAssign{v}(p)
        \in\{\True,\Undef\},\\
      \False & \text{sinon.}
    \end{cases}
  \]
  Nous montrons par induction sur~$!A$
  que (a) si
  $\pValue{v}(!A)[\LogKs]=\False$,
  alors
  $\pValue{v'}(!A)[\LogCL]=\False$,
  et que (b) si
  $\pValue{v}(!A)[\LogKs]=\True$,
  alors
  $\pValue{v'}(!A)[\LogCL]=\True$.
  \begin{enumerate}
    \item Cas de base : $!A\ident p$.
    D'après la \olref[syn][val]{defn:pValue},
    $\pValue{v}(p)[\LogKs]
    =\pAssign{v}(p)$ et
    $\pValue{v'}(p)[\LogCL]
    =\pAssign{v'}(p)$. Si la première
    valeur est $\False$, la seconde
    l'est aussi par définition de
    $\pAssign{v'}$; si elle est
    $\True$, la seconde l'est aussi.
    Ainsi (a) et (b) sont établies.

    Pour l'hérédité, examinons les cas
    suivants.
    \item $!A\ident\lnot !B$.
    \begin{enumerate}
      \item Supposons que
      $\pValue{v}(\lnot !B)[\LogKs]
      =\False$. D'après la définition
      de $\tf{\lnot}[\LogKs]$,
      $\pValue{v}(!B)[\LogKs]=\True$.
      L'hypothèse d'induction~(b)
      donne
      $\pValue{v'}(!B)[\LogCL]=\True$,
      donc
      $\pValue{v'}(\lnot !B)[\LogCL]
      =\False$.
      \item Supposons que
      $\pValue{v}(\lnot !B)[\LogKs]
      =\True$. Alors
      $\pValue{v}(!B)[\LogKs]=\False$.
      L'hypothèse d'induction~(a)
      donne
      $\pValue{v'}(!B)[\LogCL]=\False$,
      donc
      $\pValue{v'}(\lnot !B)[\LogCL]
      =\True$.
    \end{enumerate}

    \item $!A\ident(!B\land !C)$.
    \begin{enumerate}
      \item Supposons que
      $\pValue{v}(!B\land !C)[\LogKs]
      =\False$. Alors
      $\pValue{v}(!B)[\LogKs]=\False$
      ou
      $\pValue{v}(!C)[\LogKs]=\False$.
      L'hypothèse d'induction~(a)
      donne, respectivement,
      $\pValue{v'}(!B)[\LogCL]=\False$
      ou
      $\pValue{v'}(!C)[\LogCL]=\False$;
      ainsi
      $\pValue{v'}(!B\land !C)[\LogCL]
      =\False$.
      \item Supposons que
      $\pValue{v}(!B\land !C)[\LogKs]
      =\True$. Alors
      $\pValue{v}(!B)[\LogKs]=\True$
      et
      $\pValue{v}(!C)[\LogKs]=\True$.
      L'hypothèse d'induction~(b)
      donne
      $\pValue{v'}(!B)[\LogCL]=\True$
      et
      $\pValue{v'}(!C)[\LogCL]=\True$;
      ainsi
      $\pValue{v'}(!B\land !C)[\LogCL]
      =\True$.
    \end{enumerate}
  \end{enumerate}

  Les deux autres cas se traitent de
  manière analogue et sont laissés
  en exercice. On peut aussi remarquer
  que la démonstration précédente
  établit le résultat pour les formules
  ne contenant que $\lnot$ et~$\land$.
  Dans $\LogKs$ comme dans $\LogCL$,
  et pour toute valuation~$\pAssign{v}$,
  $\pValue{v}(!B\lor !C)
  =\pValue{v}(\lnot(\lnot !B\land
  \lnot !C))$ et
  $\pValue{v}(!B\lif !C)
  =\pValue{v}(\lnot(!B\land
  \lnot !C))$.
\end{proof}

\begin{prob}
  Compléter la preuve de la
  \olref[mvl][thr][mul]{prop:LP-taut-CL},
  c'est-à-dire établir (a) et (b)
  pour les cas
  $!A\ident(!B\lor !C)$ et
  $!A\ident(!B\lif !C)$.
\end{prob}

\begin{prob}
Montrer que toute tautologie classique
est une tautologie de~$\LogHal$.
\end{prob}

Bien qu'elles aient les mêmes
tautologies que la logique classique
sur le langage commun des quatre
connecteurs, leurs relations de
conséquence diffèrent. Ainsi,
$\LogLP$ est \emph{paraconsistante} :
$\lnot p,p\Entails/ q$. Le principe
d'explosion
$\lnot !A,!A\Entails !B$
ne vaut donc pas en général. Il vaut
dans certains cas, par exemple si
$!B$ est une tautologie.

\begin{prob}
  Parmi les relations suivantes,
  lesquelles valent (a) dans~$\LogLP$
  et (b) dans~$\LogHal$? Donner une
  table de vérité pour chacune.
  \begin{enumerate}
    \item $p,p\lif q\Entails q$;
    \item $\lnot q,p\lif q\Entails\lnot p$;
    \item $p\lor q,\lnot p\Entails q$;
    \item $\lnot p,p\Entails q$;
    \item $p\Entails p\lor q$;
    \item $p\lif q,q\lif r\Entails p\lif r$.
  \end{enumerate}
\end{prob}

Que se passe-t-il si l'on désigne
également $\Undef$ dans~$\LogLuk[3]$?

\begin{defn}
  La \emph{variante de \L ukasiewicz à deux valeurs
  désignées}~$\LogLuk[3]^{+}$ est définie
  par la matrice suivante :
  \begin{enumerate}
    \item Le langage propositionnel usuel
    $\Lang L_0$ avec $\lfalse$,
    $\lnot$, $\land$, $\lor$ et $\lif$.
    \item L'ensemble des valeurs de vérité
    $V=\{\True,\Undef,\False\}$.
    \item $\True$ et $\Undef$ sont
    désignées, c'est-à-dire
    $V^+=\{\True,\Undef\}$.
    \item On prend
    $\tf{\lfalse}=\False$; les
    fonctions de vérité des autres
    connecteurs sont celles de la
    logique de \L ukasiewicz~$\LogLuk[3]$.
  \end{enumerate}
\end{defn}

\begin{prob}
  Parmi les relations suivantes,
  lesquelles valent dans~$\LogLuk[3]^{+}$?
  \begin{enumerate}
    \item $p,p\lif q\Entails q$;
    \item $p\lor q,\lnot p\Entails q$;
    \item $\lnot p,p\Entails q$;
    \item $p\Entails p\lor q$.
  \end{enumerate}
\end{prob}

Des tables de vérité différentes
peuvent parfois définir la même
logique, c'est-à-dire la même relation
de conséquence, selon le choix des
valeurs désignées. C'est par exemple
le cas si, dans la logique de G\"odel,
on prend
$V^+=\{\True,\Undef\}$
au lieu de~$\{\True\}$.

\begin{prop}\ollabel{prop:gl-udes}
  La matrice ayant
  $V=\{\False,\Undef,\True\}$,
  $V^+=\{\True,\Undef\}$ et les
  fonctions de vérité de la logique
  de G\"odel à trois valeurs définit
  la logique classique.
\end{prop}

\begin{proof}
  Exercice.
\end{proof}

\begin{prob}
  Démontrer la \olref[mvl][thr][mul]{prop:gl-udes}
  en montrant que, pour la logique
  $\Log L$ définie comme celle de
  G\"odel mais avec
  $V^+=\{\True,\Undef\}$, si
  $\Gamma\Entails/[\Log L] !B$,
  alors
  $\Gamma\Entails/[\LogCL] !B$.
  S'inspirer de la
  \olref[mvl][thr][mul]{prop:LP-taut-CL},
  mais montrer, au lieu des propriétés
  (a) et (b), que
  $\pValue{v}(!A)[\LogGod]=\False$
  si et seulement si
  $\pValue{v'}(!A)[\LogCL]=\False$;
  il s'ensuit que
  $\pValue{v}(!A)[\LogGod]
  \in\{\True,\Undef\}$ si et
  seulement si
  $\pValue{v'}(!A)[\LogCL]=\True$.
  Expliquer pourquoi cela établit
  la proposition.
\end{prob}

\begin{editorial}
La preuve de la proposition sur
$\LogLP$ est formulée sur le langage
commun des quatre connecteurs :
l'opérateur~$+$ de~$\LogHal$ ne
fait pas partie du langage classique
ainsi comparé. Dans le cas de base
de l'induction, la source affirme
à tort l'égalité des valeurs des
variables sous $\pAssign{v}$ et
$\pAssign{v'}$ même lorsque la
première vaut $\Undef$; seules les
deux implications conditionnelles
(a) et (b) sont nécessaires.
Dans le cas de la conjonction,
la source répète $!B$ là où elle
doit aussi considérer~$!C$.
La source appelle $\LogRM[3]$ cette variante,
notée ici $\LogLuk[3]^{+}$; nous conservons ses tables
et complétons sa constante absente de $\LogLuk[3]$ par $\tf{\lfalse}=\False$.
Elle diffère de la logique usuelle R-Mingle~3 :
pour $(\True,\Undef)$, $(\Undef,\False)$ et $(\Undef,\Undef)$,
son conditionnel donne $\Undef$, $\Undef$ et $\True$ au lieu de $\False$, $\False$ et $\Undef$.
Voir \href{https://cdn.aaai.org/ocs/15415/15415-68643-1-PB.pdf}{Sutcliffe, Pelletier et Hazen (2017), table~1}.
\end{editorial}

\end{document}
